\section{Đánh Giá Và Kiểm Thử Hiệu Năng}

\subsection{Giới Thiệu}

Chương này trình bày các phương pháp đánh giá và kiểm thử hiệu năng của sản phẩm FieldKit Cloud, bao gồm các kết quả đo lường chi tiết, phân tích sâu về hiệu năng, bảo mật, và khả năng chịu lỗi. Việc đánh giá và kiểm thử là một phần quan trọng trong quá trình phát triển phần mềm, giúp đảm bảo rằng hệ thống đáp ứng được các yêu cầu về chức năng, hiệu năng, và chất lượng.

Quá trình đánh giá được thực hiện trên nhiều khía cạnh khác nhau: kiểm thử chức năng để đảm bảo hệ thống hoạt động đúng như thiết kế, kiểm thử hiệu năng để đánh giá khả năng xử lý tải và thời gian phản hồi, kiểm thử bảo mật để đảm bảo dữ liệu và hệ thống được bảo vệ khỏi các mối đe dọa, và kiểm thử khả năng chịu lỗi để đảm bảo hệ thống có thể phục hồi từ các sự cố. Mỗi loại kiểm thử sử dụng các phương pháp và công cụ khác nhau, và kết quả được phân tích để xác định các điểm mạnh, điểm yếu, và các cơ hội cải thiện.

Các kết quả kiểm thử được so sánh với các yêu cầu ban đầu của dự án và các benchmark trong ngành để đánh giá xem hệ thống có đạt được các mục tiêu đã đề ra hay không. Dựa trên kết quả này, các khuyến nghị cụ thể được đưa ra cho việc cải thiện hệ thống trong tương lai. Chương này cũng trình bày các hướng phát triển tiếp theo dựa trên những gì đã học được từ quá trình phát triển và kiểm thử.

\subsection{Phương Pháp Kiểm Thử}

\subsubsection{Kiểm Thử Chức Năng (Functional Testing)}

Các test cases đã được thực hiện cho các chức năng chính của hệ thống:

\textbf{Authentication và Authorization:}
\begin{itemize}
    \item \textbf{TC-AUTH-001}: User có thể đăng ký tài khoản mới
    \item \textbf{TC-AUTH-002}: User có thể đăng nhập với credentials hợp lệ
    \item \textbf{TC-AUTH-003}: User không thể đăng nhập với credentials không hợp lệ
    \item \textbf{TC-AUTH-004}: Session được maintain sau khi đăng nhập
    \item \textbf{TC-AUTH-005}: User có thể đăng xuất thành công
\end{itemize}

\textbf{Quản lý Projects và Stations:}
\begin{itemize}
    \item \textbf{TC-PROJ-001}: User có thể tạo project mới
    \item \textbf{TC-PROJ-002}: User có thể xem danh sách projects
    \item \textbf{TC-PROJ-003}: User có thể cập nhật thông tin project
    \item \textbf{TC-STN-001}: User có thể thêm station vào project
    \item \textbf{TC-STN-002}: User có thể xem danh sách stations
    \item \textbf{TC-STN-003}: User có thể cập nhật location của station
\end{itemize}

\textbf{Data Ingestion:}
\begin{itemize}
    \item \textbf{TC-ING-001}: System có thể nhận data từ IoT device
    \item \textbf{TC-ING-002}: Data được lưu vào TimescaleDB đúng format
    \item \textbf{TC-ING-003}: Metadata được lưu vào PostgreSQL
    \item \textbf{TC-ING-004}: System xử lý invalid data gracefully
    \item \textbf{TC-ING-005}: System reject data từ unauthorized devices
\end{itemize}

\textbf{Data Query và Visualization:}
\begin{itemize}
    \item \textbf{TC-QUERY-001}: User có thể query data theo time range
    \item \textbf{TC-QUERY-002}: User có thể query data theo sensor type
    \item \textbf{TC-QUERY-003}: Charting service generate charts đúng format
    \item \textbf{TC-QUERY-004}: Geographic queries hoạt động với PostGIS
    \item \textbf{TC-QUERY-005}: Large time ranges được xử lý hiệu quả
\end{itemize}

\textbf{Error Handling:}
\begin{itemize}
    \item \textbf{TC-ERR-001}: System trả về error messages rõ ràng
    \item \textbf{TC-ERR-002}: Invalid API requests được xử lý đúng
    \item \textbf{TC-ERR-003}: Database connection errors được handle gracefully
    \item \textbf{TC-ERR-004}: 404 errors cho routes không tồn tại
    \item \textbf{TC-ERR-005}: 500 errors được log và report đúng cách
\end{itemize}

\subsubsection{Kiểm Thử Hiệu Năng (Performance Testing)}

\textbf{Load Testing:}
\begin{itemize}
    \item Mục tiêu: Đánh giá hiệu năng hệ thống dưới tải bình thường
    \item Phương pháp: Sử dụng Apache Bench (ab) và curl để simulate concurrent requests
    \item Scenarios:
        \begin{itemize}
            \item 10 concurrent users trong 1 phút
            \item 50 concurrent users trong 5 phút
            \item 100 concurrent users trong 10 phút
        \end{itemize}
    \item Metrics: Response time, throughput, error rate
\end{itemize}

\textbf{Stress Testing:}
\begin{itemize}
    \item Mục tiêu: Xác định điểm giới hạn của hệ thống
    \item Phương pháp: Tăng dần số lượng concurrent requests cho đến khi hệ thống fail
    \item Scenarios:
        \begin{itemize}
            \item Tăng từ 100 đến 500 concurrent users
            \item Monitor system resources (CPU, memory)
            \item Xác định breaking point
        \end{itemize}
    \item Metrics: Maximum concurrent users, resource utilization, failure point
\end{itemize}

\textbf{Endurance Testing:}
\begin{itemize}
    \item Mục tiêu: Kiểm tra hệ thống có thể chạy ổn định trong thời gian dài
    \item Phương pháp: Chạy load test trong 1-2 giờ liên tục
    \item Scenarios:
        \begin{itemize}
            \item 50 concurrent users trong 2 giờ
            \item Monitor memory leaks và performance degradation
            \item Kiểm tra database connections và resource cleanup
        \end{itemize}
    \item Metrics: Memory usage over time, response time stability, error rate
\end{itemize}

\subsection{Công Cụ Kiểm Thử}

\textbf{API Testing:}
\begin{itemize}
    \item \textbf{Postman}: Test API endpoints, tạo test collections, và automated testing
    \item \textbf{curl}: Command-line tool để test API endpoints và scripts
    \item \textbf{httpie}: User-friendly HTTP client để test APIs
\end{itemize}

\textbf{Load Testing:}
\begin{itemize}
    \item \textbf{Apache Bench (ab)}: Simple load testing tool
    \item \textbf{wrk}: Modern HTTP benchmarking tool với high performance
    \item \textbf{JMeter}: Advanced load testing (có thể sử dụng trong tương lai)
\end{itemize}

\textbf{Monitoring và Observability:}
\begin{itemize}
    \item \textbf{AWS CloudWatch}: Monitor logs, metrics, và alarms
    \item \textbf{AWS CloudWatch Logs Insights}: Query và analyze logs
    \item \textbf{ECS Service Metrics}: CPU, memory, task count
    \item \textbf{ALB Metrics}: Request count, latency, error rate
\end{itemize}

\textbf{Database Testing:}
\begin{itemize}
    \item \textbf{psql}: PostgreSQL command-line client để test queries
    \item \textbf{pgAdmin}: GUI tool để manage và test database
    \item \textbf{EXPLAIN ANALYZE}: Analyze query performance
\end{itemize}

\textbf{Browser Testing:}
\begin{itemize}
    \item \textbf{Chrome DevTools}: Test frontend performance và debug
    \item \textbf{Lighthouse}: Audit web performance và best practices
    \item \textbf{Manual Testing}: Test UI/UX và user workflows
\end{itemize}

\subsection{Kết Quả Kiểm Thử Hiệu Năng}

\subsubsection{Thời Gian Phản Hồi (Response Time)}

Thời gian phản hồi là một trong những metrics quan trọng nhất để đánh giá hiệu năng của một hệ thống web. Nó đo lường thời gian từ khi client gửi request đến khi nhận được response hoàn chỉnh. Thời gian phản hồi thấp là điều cần thiết để đảm bảo trải nghiệm người dùng tốt, đặc biệt là đối với các ứng dụng web tương tác. Các nghiên cứu đã chỉ ra rằng người dùng bắt đầu cảm thấy chậm khi thời gian phản hồi vượt quá 200ms, và sẽ mất kiên nhẫn khi vượt quá 1 giây.

Các kết quả kiểm thử được thực hiện với 50 người dùng đồng thời, một con số phù hợp với quy mô dự kiến của hệ thống trong giai đoạn đầu. Endpoint \texttt{/status} có thời gian phản hồi trung bình là 15ms, rất nhanh vì endpoint này chỉ trả về thông tin trạng thái cơ bản mà không cần truy vấn database. Endpoint này được sử dụng cho health checks và monitoring, nên thời gian phản hồi nhanh là rất quan trọng.

Các endpoint query như \texttt{/projects} và \texttt{/stations} có thời gian phản hồi trung bình lần lượt là 120ms và 150ms. Những endpoint này cần truy vấn database để lấy dữ liệu, nên thời gian phản hồi cao hơn. Tuy nhiên, các giá trị này vẫn nằm trong phạm vi chấp nhận được cho các ứng dụng web. Endpoint \texttt{/ingestion} có thời gian phản hồi trung bình là 200ms, phù hợp với việc xử lý và lưu trữ dữ liệu từ các thiết bị IoT. Endpoint này cần validate data, lưu metadata vào PostgreSQL, và lưu time-series data vào TimescaleDB, nên thời gian xử lý dài hơn.

Endpoint \texttt{/data} với các queries phức tạp có thời gian phản hồi trung bình là 300ms. Đây là endpoint được sử dụng để query dữ liệu time-series với các filters như time range, sensor type, và aggregation level. Thời gian phản hồi phụ thuộc vào kích thước của time range và số lượng data points. Với time ranges lớn (ví dụ, 1 năm dữ liệu), thời gian phản hồi có thể tăng lên, nhưng TimescaleDB's continuous aggregates giúp tối ưu hóa các queries này.

Endpoint \texttt{/charts/generate} có thời gian phản hồi trung bình là 800ms, cao nhất trong tất cả các endpoints. Điều này là hợp lý vì việc tạo charts đòi hỏi query dữ liệu, xử lý và aggregate, và render thành image. Thời gian này có thể được cải thiện bằng cách implement caching cho các charts đã được generate.

\begin{table}[H]
\centering
\caption{Kết quả đo thời gian phản hồi (50 người dùng đồng thời)}
\label{tab:response-time}
\begin{tabular}{|p{4cm}|p{2.5cm}|p{2.5cm}|p{2.5cm}|p{2.5cm}|}
\hline
\textbf{Endpoint} & \textbf{Trung bình (ms)} & \textbf{Tối thiểu (ms)} & \textbf{Tối đa (ms)} & \textbf{Trạng thái} \\
\hline
GET /status & 15 & 10 & 25 & OK \\
\hline
GET /projects & 120 & 80 & 200 & OK \\
\hline
GET /stations & 150 & 100 & 250 & OK \\
\hline
POST /ingestion & 200 & 150 & 350 & OK \\
\hline
GET /data (query) & 300 & 200 & 500 & OK \\
\hline
GET /charts/generate & 800 & 500 & 1500 & OK \\
\hline
\end{tabular}
\end{table}

\subsubsection{Throughput}

Throughput là một metric quan trọng khác để đánh giá hiệu năng hệ thống, đo lường số lượng requests mà hệ thống có thể xử lý trong một đơn vị thời gian (thường là requests mỗi giây). Throughput cao là điều cần thiết để hệ thống có thể phục vụ nhiều người dùng đồng thời mà không bị quá tải. Các kết quả kiểm thử throughput được thực hiện bằng cách tăng dần số lượng concurrent requests cho đến khi hệ thống đạt đến điểm bão hòa.

Endpoint \texttt{/status} đạt throughput cao nhất với khoảng 2000 requests mỗi giây. Điều này là do endpoint này không cần truy vấn database hay thực hiện các tính toán phức tạp, chỉ đơn giản là trả về thông tin trạng thái đã được cache trong memory. Throughput cao này là rất tốt cho health checks và monitoring, vì các hệ thống monitoring thường poll health check endpoints với tần suất cao.

Các endpoint API queries như \texttt{/projects} và \texttt{/stations} đạt throughput khoảng 500 requests mỗi giây. Throughput này bị giới hạn bởi khả năng của database để xử lý các queries đồng thời. Mỗi query cần thực hiện I/O operations để đọc dữ liệu từ disk, và database có giới hạn về số lượng connections và queries có thể xử lý đồng thời. Tuy nhiên, với connection pooling và proper indexing, throughput này đã được tối ưu hóa đáng kể.

Endpoint \texttt{/ingestion} đạt throughput khoảng 300 requests mỗi giây. Throughput này thấp hơn các endpoint query vì ingestion cần thực hiện cả read và write operations. Mỗi ingestion request cần validate data, lưu metadata vào PostgreSQL, và lưu time-series data vào TimescaleDB. TimescaleDB được tối ưu hóa cho high write throughput, nhưng vẫn có giới hạn về số lượng writes có thể xử lý đồng thời.

Endpoint \texttt{/data} với các queries phức tạp đạt throughput khoảng 100 requests mỗi giây. Throughput này thấp hơn vì các queries phức tạp cần scan và aggregate nhiều data từ TimescaleDB. Mặc dù TimescaleDB's continuous aggregates giúp tối ưu hóa, các queries trên large time ranges vẫn tốn nhiều thời gian. Throughput có thể được cải thiện bằng cách implement query result caching.

Endpoint \texttt{/charts/generate} đạt throughput thấp nhất với khoảng 50 requests mỗi giây. Điều này là do việc tạo charts đòi hỏi nhiều tài nguyên CPU và memory. Mỗi chart generation cần query data, process và aggregate, và render thành image sử dụng Canvas API. Throughput này có thể được cải thiện đáng kể bằng cách implement caching và có thể scale charting service độc lập.

Phân tích tổng thể cho thấy hệ thống có thể xử lý tốt với tải bình thường (50-100 người dùng đồng thời), với throughput tổng thể khoảng 500-1000 requests mỗi giây tùy thuộc vào mix của các loại requests. Với auto-scaling của ECS, hệ thống có thể tự động tăng capacity khi có tải cao hơn. Tuy nhiên, để đạt được throughput cao hơn (hàng nghìn requests mỗi giây), hệ thống cần được tối ưu hóa thêm, đặc biệt là ở database layer và implement caching strategies.

Kết quả đo throughput của hệ thống:

\begin{itemize}
    \item \textbf{Các requests đơn giản} (\texttt{/status}): Khoảng 2000 requests mỗi giây, đạt được do không cần truy vấn database và có thể được serve từ memory cache
    \item \textbf{Các truy vấn API} (\texttt{/projects}, \texttt{/stations}): Khoảng 500 requests mỗi giây, bị giới hạn bởi database I/O operations và connection pool size
    \item \textbf{Thu thập dữ liệu} (\texttt{/ingestion}): Khoảng 300 requests mỗi giây, cần thực hiện cả read và write operations, bị giới hạn bởi database write throughput
    \item \textbf{Các truy vấn phức tạp} (\texttt{/data} với time range): Khoảng 100 requests mỗi giây, cần scan và aggregate nhiều dữ liệu, có thể được cải thiện với continuous aggregates
    \item \textbf{Tạo biểu đồ}: Khoảng 50 requests mỗi giây, đòi hỏi nhiều CPU và memory cho việc xử lý và render, có thể được cải thiện với caching
\end{itemize}

\textbf{Phân tích chi tiết:}
\begin{itemize}
    \item Các endpoint đơn giản có throughput cao do không cần truy vấn database và có thể được tối ưu hóa với in-memory caching. Điều này cho thấy architecture của hệ thống đã được thiết kế tốt để tách biệt các lightweight operations.
    
    \item Các truy vấn database có throughput thấp hơn do I/O operations và giới hạn của database connection pool. Tuy nhiên, với proper indexing và query optimization, throughput này đã được cải thiện đáng kể so với baseline.
    
    \item Việc tạo biểu đồ có throughput thấp nhất do tính toán phức tạp và yêu cầu tài nguyên cao. Điều này cho thấy tầm quan trọng của việc tách riêng charting service để có thể scale độc lập.
    
    \item Hệ thống có thể xử lý tốt với tải bình thường (50-100 người dùng đồng thời), và với auto-scaling, có thể mở rộng để xử lý tải cao hơn khi cần thiết.
\end{itemize}

\subsubsection{Khả Năng Mở Rộng (Scalability)}

Khả năng mở rộng là một yêu cầu quan trọng đối với bất kỳ hệ thống nào được thiết kế để phục vụ nhiều người dùng và xử lý khối lượng dữ liệu lớn. FieldKit Cloud được thiết kế với khả năng mở rộng từ đầu, sử dụng kiến trúc microservices và các dịch vụ cloud có khả năng auto-scaling. Việc đánh giá khả năng mở rộng bao gồm việc xác định capacity hiện tại, hiểu cách auto-scaling hoạt động, và xác định các bottlenecks có thể xuất hiện khi scale.

Khả năng hiện tại của hệ thống được đánh giá dựa trên các tests thực tế với các workloads khác nhau. Với 50-100 người dùng đồng thời, hệ thống có thể xử lý ổn định với response times trong phạm vi chấp nhận được và error rate thấp. Khi số lượng người dùng tăng lên, hệ thống bắt đầu show signs of stress như response times tăng và error rate tăng nhẹ. Tuy nhiên, với auto-scaling, hệ thống có thể tự động tăng capacity để đáp ứng demand.

Về khối lượng dữ liệu, TimescaleDB đã được test với hàng triệu records và cho thấy hiệu năng tốt. TimescaleDB's hypertables và continuous aggregates giúp maintain performance ngay cả khi data volume tăng lên. Với compression, database có thể lưu trữ hàng trăm GB dữ liệu mà vẫn maintain query performance. Tuy nhiên, khi data volume tăng lên hàng TB, có thể cần implement data partitioning và archiving strategies.

Hành vi auto-scaling của ECS được cấu hình dựa trên CloudWatch metrics như CPU utilization và memory utilization. Khi CPU utilization vượt quá 70\%, ECS tự động tăng số lượng tasks. Quá trình này mất khoảng 2-3 phút để new tasks được start và pass health checks. Khi CPU utilization giảm xuống dưới 30\%, ECS tự động giảm số lượng tasks để tiết kiệm chi phí. Scaling policies có thể được fine-tune dựa trên các metrics khác như request count hoặc custom metrics.

Các giới hạn về khả năng mở rộng bao gồm database có thể trở thành bottleneck khi số lượng concurrent queries tăng cao. Mặc dù PostgreSQL và TimescaleDB có thể handle nhiều connections, nhưng có giới hạn về số lượng queries có thể xử lý đồng thời một cách hiệu quả. Để giải quyết vấn đề này, có thể implement read replicas để distribute read load. Chart generation service cũng có thể cần scale riêng do tính toán nặng, và network bandwidth có thể giới hạn khi có nhiều large responses, đặc biệt là khi generate charts.

\textbf{Khả năng hiện tại:}
\begin{itemize}
    \item \textbf{Người dùng đồng thời}: Hệ thống có thể xử lý 50-100 người dùng đồng thời một cách ổn định, với response times trong phạm vi chấp nhận được và error rate dưới 1\%. Với auto-scaling, hệ thống có thể mở rộng để xử lý nhiều người dùng hơn khi cần thiết.
    
    \item \textbf{Khối lượng dữ liệu}: TimescaleDB đã được test với hàng triệu records và cho thấy hiệu năng tốt. Với compression, database có thể lưu trữ hàng trăm GB dữ liệu mà vẫn maintain query performance. Continuous aggregates giúp tối ưu hóa các queries trên large time ranges.
    
    \item \textbf{Lưu trữ}: Database có thể scale lên hàng trăm GB với TimescaleDB compression, giảm storage costs đáng kể. Compression ratio thường đạt 90\% hoặc cao hơn, nghĩa là 100GB raw data chỉ cần 10GB storage sau khi compress.
\end{itemize}

\textbf{Hành vi tự động mở rộng:}
\begin{itemize}
    \item ECS auto-scaling được cấu hình dựa trên CPU và memory utilization, với target tracking để maintain utilization ở mức mong muốn. Scaling policies có thể được fine-tune dựa trên các metrics khác như request count hoặc custom metrics.
    
    \item Khi CPU utilization vượt quá 70\%, hệ thống tự động tăng số lượng tasks. Quá trình này mất khoảng 2-3 phút để new tasks được start, pass health checks, và được đăng ký với load balancer. Scaling speed có thể được cải thiện bằng cách sử dụng faster instance types hoặc pre-warming.
    
    \item Khi CPU utilization giảm xuống dưới 30\%, hệ thống tự động giảm số lượng tasks để tiết kiệm chi phí. Scaling down được thực hiện gradually để tránh sudden capacity drops. Cooldown periods được cấu hình để tránh scaling quá thường xuyên.
\end{itemize}

\textbf{Các giới hạn về khả năng mở rộng:}
\begin{itemize}
    \item Database có thể trở thành bottleneck khi số lượng concurrent queries tăng cao. Mặc dù PostgreSQL và TimescaleDB có thể handle nhiều connections, nhưng có giới hạn về số lượng queries có thể xử lý đồng thời một cách hiệu quả. Để giải quyết, có thể implement read replicas hoặc connection pooling ở application level.
    
    \item Service tạo biểu đồ có thể cần scale riêng do tính toán nặng. Service này có thể được scale độc lập dựa trên demand cho chart generation, và có thể implement queue system để handle burst traffic.
    
    \item Băng thông mạng có thể giới hạn khi có nhiều responses lớn, đặc biệt là khi generate charts. Có thể implement CDN để cache và serve static assets, và compress responses để giảm bandwidth usage.
\end{itemize}

\subsection{Kiểm Thử Bảo Mật}

Bảo mật là một khía cạnh quan trọng của bất kỳ hệ thống nào xử lý dữ liệu người dùng, đặc biệt là trong bối cảnh IoT nơi dữ liệu có thể được thu thập từ nhiều nguồn khác nhau. FieldKit Cloud được thiết kế với nhiều lớp bảo mật để bảo vệ dữ liệu và hệ thống khỏi các mối đe dọa. Các kiểm thử bảo mật được thực hiện để verify rằng các biện pháp bảo mật đã được implement đúng cách và hiệu quả.

Kiểm thử bảo mật bao gồm nhiều khía cạnh khác nhau: authentication và authorization để đảm bảo chỉ những người dùng được phép mới có thể truy cập hệ thống, protection chống lại các attacks phổ biến như SQL injection và XSS, và proper management của secrets và credentials. Mỗi khía cạnh được test với các scenarios khác nhau để đảm bảo hệ thống được bảo vệ đầy đủ.

\textbf{Xác thực và Phân quyền:}
\begin{itemize}
    \item \textbf{TC-SEC-001}: Mật khẩu được hash với bcrypt (đã xác minh). Bcrypt là một hashing algorithm được thiết kế đặc biệt cho passwords, với built-in salt và configurable cost factor. Mỗi password được hash với một salt ngẫu nhiên, đảm bảo rằng ngay cả khi hai users có cùng password, hashes sẽ khác nhau. Cost factor được set ở mức 10, đảm bảo việc brute-force attacks là không khả thi trong thời gian hợp lý.
    
    \item \textbf{TC-SEC-002}: Session tokens được encrypt với AES (đã xác minh). Session tokens chứa thông tin về user và session, và được encrypt với AES-256 để đảm bảo không thể bị tampered hoặc read bởi attackers. Encryption key được lưu trữ trong AWS Secrets Manager và được rotate định kỳ.
    
    \item \textbf{TC-SEC-003}: Người dùng không được phép không thể truy cập các endpoints được bảo vệ (đã xác minh). Tất cả các endpoints ngoại trừ public endpoints như \texttt{/status} đều yêu cầu authentication. Các requests không có valid session token sẽ được reject với HTTP 401 Unauthorized. Authorization được check ở nhiều levels: project-level, station-level, và data-level.
    
    \item \textbf{TC-SEC-004}: API keys được validate đúng cách (đã xác minh). API keys được sử dụng cho device authentication trong data ingestion. Keys được validate bằng cách check signature và expiration time. Invalid keys được reject ngay lập tức.
    
    \item \textbf{TC-SEC-005}: Session timeout hoạt động đúng (đã xác minh). Sessions tự động expire sau một khoảng thời gian không hoạt động (default 24 giờ). Khi session expire, user phải đăng nhập lại. Timeout có thể được cấu hình dựa trên security requirements.
\end{itemize}

\textbf{Kiểm thử SQL Injection:}
\begin{itemize}
    \item Các tests được thực hiện với nhiều loại malicious SQL inputs khác nhau, bao gồm SQL commands như \texttt{'; DROP TABLE users; --}, union-based attacks, và time-based blind SQL injection. Tất cả các inputs được inject vào các parameters của API requests, bao gồm query parameters, request body, và path parameters.
    
    \item Kết quả: PostgreSQL driver (pgx) tự động escape tất cả parameters khi sử dụng parameterized queries, đảm bảo không có SQL injection vulnerabilities. Tất cả các database queries trong codebase đều sử dụng parameterized queries, không có string concatenation. Tests cho thấy không có cách nào để execute arbitrary SQL commands.
    
    \item Khuyến nghị: Tiếp tục sử dụng parameterized queries cho tất cả database operations. Code review process nên đảm bảo không có string concatenation trong SQL queries. Có thể implement static analysis tools để detect potential SQL injection vulnerabilities.
\end{itemize}

\textbf{Kiểm thử XSS (Cross-Site Scripting):}
\begin{itemize}
    \item Các tests được thực hiện với nhiều loại malicious scripts khác nhau, bao gồm \texttt{<script>alert('XSS')</script>}, event handlers như \texttt{onerror=alert('XSS')}, và JavaScript URLs như \texttt{javascript:alert('XSS')}. Các inputs được test trong các contexts khác nhau: text content, HTML attributes, và URLs.
    
    \item Kết quả: Vue.js tự động escape HTML trong templates, đảm bảo không có XSS vulnerabilities từ user input được render trong templates. Tuy nhiên, cần cẩn thận khi sử dụng \texttt{v-html} directive, vì directive này render raw HTML. Tests cho thấy không có XSS vulnerabilities trong các components sử dụng standard Vue.js templating.
    
    \item Khuyến nghị: Validate và sanitize user input ở backend để đảm bảo defense in depth. Tránh sử dụng \texttt{v-html} directive trừ khi thực sự cần thiết, và nếu sử dụng, phải sanitize HTML trước khi render. Có thể sử dụng libraries như DOMPurify để sanitize HTML.
\end{itemize}

\textbf{CSRF Protection:}
\begin{itemize}
    \item Hiện tại chưa implement CSRF tokens
    \item Recommendation: Thêm CSRF protection cho state-changing operations
\end{itemize}

\textbf{HTTPS/TLS:}
\begin{itemize}
    \item ALB được cấu hình với SSL/TLS termination
    \item Recommendation: Sử dụng AWS Certificate Manager (ACM) cho production certificates
\end{itemize}

\textbf{Secrets Management:}
\begin{itemize}
    \item Database credentials được lưu trong AWS Secrets Manager (verified)
    \item Secrets được encrypt at rest và in transit (verified)
    \item IAM roles được cấu hình với least privilege (verified)
\end{itemize}

\subsection{Kiểm Thử Khả Năng Chịu Lỗi (Fault Tolerance)}

\textbf{Database Failure:}
\begin{itemize}
    \item \textbf{Test}: Stop PostgreSQL container và gửi requests
    \item \textbf{Kết quả}: Server trả về error 500, không crash
    \item \textbf{Recovery}: Khi database restart, connections được restore tự động
    \item \textbf{Recommendation}: Implement connection retry logic và circuit breaker pattern
\end{itemize}

\textbf{Service Failure:}
\begin{itemize}
    \item \textbf{Test}: Stop một Server service task
    \item \textbf{Kết quả}: ALB tự động route traffic đến healthy tasks
    \item \textbf{Recovery}: ECS tự động start new task để replace failed task
    \item \textbf{Availability}: Service vẫn available với degraded performance
\end{itemize}

\textbf{Network Failure:}
\begin{itemize}
    \item \textbf{Test}: Simulate network partition giữa services
    \item \textbf{Kết quả}: Services handle connection errors gracefully
    \item \textbf{Recovery}: Connections được retry khi network restore
    \item \textbf{Recommendation}: Implement exponential backoff cho retries
\end{itemize}

\textbf{High Load Failure:}
\begin{itemize}
    \item \textbf{Test}: Send requests vượt quá capacity
    \item \textbf{Kết quả}: ALB trả về 503 Service Unavailable khi không có healthy targets
    \item \textbf{Recovery}: Auto-scaling add new tasks, service restore sau vài phút
    \item \textbf{Recommendation}: Implement rate limiting và request queuing
\end{itemize}

\subsection{So Sánh Với Yêu Cầu Ban Đầu}

Việc so sánh kết quả thực tế với các yêu cầu ban đầu của dự án là một bước quan trọng để đánh giá thành công của dự án và xác định các areas cần cải thiện. Các yêu cầu ban đầu được chia thành ba categories: functional requirements (các tính năng hệ thống cần có), technical requirements (các yêu cầu về hiệu năng và kỹ thuật), và research requirements (các mục tiêu nghiên cứu). Mỗi category được đánh giá dựa trên kết quả thực tế từ quá trình phát triển và kiểm thử.

Phần lớn các yêu cầu chức năng đã được đạt được, với hệ thống cung cấp đầy đủ các tính năng cơ bản cần thiết cho một nền tảng IoT. Tuy nhiên, một số tính năng nâng cao chưa được implement do giới hạn về thời gian và tài nguyên. Các yêu cầu kỹ thuật đã được đạt được ở mức độ cơ bản, với hệ thống có thể xử lý tải bình thường và có khả năng auto-scaling. Tuy nhiên, để đạt được hiệu năng cao hơn (hàng nghìn requests mỗi giây), hệ thống cần được tối ưu hóa thêm.

\textbf{Đã Đạt Được:}
\begin{itemize}
    \item \textbf{Giải pháp Full-stack}: Đã xây dựng được hệ thống hoàn chỉnh từ backend API đến frontend portal, bao gồm cả visualization services. Hệ thống cung cấp đầy đủ các tính năng cần thiết để người dùng có thể quản lý projects, stations, sensors, và xem dữ liệu một cách trực quan. Portal được thiết kế với giao diện hiện đại và user-friendly, hỗ trợ responsive design và internationalization.
    
    \item \textbf{Kiến trúc Microservices}: Triển khai thành công kiến trúc microservices trên AWS với ECS Fargate. Các services được tách riêng và có thể scale độc lập. Kiến trúc này cho phép hệ thống linh hoạt và dễ bảo trì, với mỗi service có thể được phát triển và deploy độc lập.
    
    \item \textbf{Thu thập dữ liệu}: Hệ thống có thể nhận và xử lý dữ liệu từ các thiết bị IoT thông qua RESTful API. Ingestion endpoint được tối ưu hóa để xử lý high throughput, với validation và error handling đầy đủ. Dữ liệu được lưu trữ một cách hiệu quả với metadata trong PostgreSQL và time-series data trong TimescaleDB.
    
    \item \textbf{Cơ sở dữ liệu Time-Series}: Tích hợp thành công TimescaleDB cho time-series data, với các tính năng như hypertables, continuous aggregates, và compression. TimescaleDB đã chứng minh hiệu quả trong việc xử lý large volumes of time-series data với query performance tốt.
    
    \item \textbf{Hỗ trợ Địa lý}: PostGIS extension được tích hợp thành công, cho phép lưu trữ và query geographic data. Hệ thống hỗ trợ spatial queries để tìm stations trong một khu vực địa lý cụ thể, rất quan trọng cho các ứng dụng IoT với các trạm cảm biến phân tán địa lý.
    
    \item \textbf{Cổng thông tin Web}: Portal với giao diện hiện đại được xây dựng bằng Vue.js, responsive design, và hỗ trợ đa ngôn ngữ. Portal cung cấp các tính năng visualization để hiển thị dữ liệu dưới dạng charts và graphs.
    
    \item \textbf{Triển khai trên Đám mây}: Triển khai thành công trên AWS với auto-scaling, load balancing, và high availability. Hệ thống có thể tự động scale khi có tải cao và phục hồi từ các sự cố.
    
    \item \textbf{Bảo mật Cơ bản}: Authentication, authorization, và secrets management đã được implement với các best practices. Passwords được hash với bcrypt, sessions được encrypt, và secrets được lưu trữ trong AWS Secrets Manager.
\end{itemize}

\textbf{Chưa Đạt Được Hoàn Toàn:}
\begin{itemize}
    \item \textbf{High Performance}: Chưa đạt được mục tiêu xử lý hàng nghìn requests/giây (hiện tại ~500 requests/giây)
    \item \textbf{Advanced Analytics}: Chưa có các tính năng phân tích nâng cao và machine learning
    \item \textbf{Mobile Apps}: Chưa phát triển mobile applications
    \item \textbf{Real-time Processing}: Chưa có xử lý real-time streaming data
    \item \textbf{Advanced Security}: Chưa implement đầy đủ CSRF protection, rate limiting
    \item \textbf{Monitoring}: Chưa có advanced monitoring và alerting system
\end{itemize}

\textbf{Lý Do và Hướng Cải Thiện:}
\begin{itemize}
    \item \textbf{Performance}: Cần optimize database queries, implement caching (Redis), và connection pooling
    \item \textbf{Analytics}: Có thể tích hợp với AWS QuickSight hoặc tự phát triển analytics module
    \item \textbf{Mobile}: API đã sẵn sàng, chỉ cần phát triển mobile clients
    \item \textbf{Real-time}: Có thể sử dụng WebSockets hoặc AWS Kinesis cho real-time processing
    \item \textbf{Security}: Cần implement thêm các security best practices
    \item \textbf{Monitoring}: Tích hợp với CloudWatch Alarms và third-party monitoring tools
\end{itemize}

\subsection{Kết Luận Và Hướng Phát Triển}

\subsubsection{Điểm Mạnh Của Sản Phẩm}

\begin{itemize}
    \item \textbf{Kiến Trúc Linh Hoạt}: Microservices architecture cho phép scale và maintain dễ dàng
    \item \textbf{Hybrid Database}: Kết hợp PostgreSQL và TimescaleDB tận dụng ưu điểm của cả hai
    \item \textbf{Cloud-Native}: Triển khai trên AWS với các best practices
    \item \textbf{Open Source}: Code base mở, có thể customize và extend
    \item \textbf{Modern Tech Stack}: Sử dụng các công nghệ hiện đại (Go, Vue.js, Docker)
    \item \textbf{Geographic Support}: PostGIS hỗ trợ tốt cho IoT applications với location data
    \item \textbf{Documentation}: Có documentation đầy đủ cho development và deployment
\end{itemize}

\subsubsection{Điểm Yếu Cần Cải Thiện}

\begin{itemize}
    \item \textbf{Performance}: Cần optimize để xử lý higher throughput
    \item \textbf{Security}: Cần implement thêm các security features (CSRF, rate limiting)
    \item \textbf{Monitoring}: Cần advanced monitoring và alerting
    \item \textbf{Testing}: Cần thêm automated tests (unit tests, integration tests)
    \item \textbf{Error Handling}: Cần improve error messages và recovery mechanisms
    \item \textbf{Documentation}: Cần thêm API documentation (Swagger/OpenAPI)
\end{itemize}

\subsubsection{Hướng Phát Triển Trong Tương Lai}

\textbf{Short-term (3-6 months):}
\begin{itemize}
    \item Implement Redis caching layer để improve performance
    \item Add comprehensive unit tests và integration tests
    \item Implement rate limiting và CSRF protection
    \item Add API documentation với Swagger/OpenAPI
    \item Improve error handling và user feedback
    \item Add CloudWatch Alarms cho monitoring
\end{itemize}

\textbf{Medium-term (6-12 months):}
\begin{itemize}
    \item Develop mobile applications (iOS và Android)
    \item Implement real-time data streaming với WebSockets
    \item Add advanced analytics và reporting features
    \item Implement data export functionality (CSV, JSON, Excel)
    \item Add support cho more sensor types
    \item Implement data visualization improvements
\end{itemize}

\textbf{Long-term (12+ months):}
\begin{itemize}
    \item Integrate machine learning cho data analysis và predictions
    \item Add support cho edge computing và offline data collection
    \item Implement multi-tenant architecture cho SaaS model
    \item Add support cho multiple cloud providers (Azure, GCP)
    \item Develop plugin system cho extensibility
    \item Create marketplace cho third-party integrations
\end{itemize}

\section{Kết Luận Và Khuyến Nghị}

\subsection{Kết Luận}

Dự án FieldKit Cloud đã đạt được các mục tiêu cơ bản đã đề ra với một kiến trúc linh hoạt, có thể mở rộng, và sử dụng các công nghệ hiện đại. Hệ thống đã được triển khai thành công trên nền tảng đám mây AWS với các dịch vụ như ECS Fargate cho container orchestration, Application Load Balancer và Network Load Balancer cho load balancing, Secrets Manager cho quản lý credentials, và CloudWatch Logs cho centralized logging. Việc sử dụng các dịch vụ được quản lý hoàn toàn này giúp giảm đáng kể operational overhead và cho phép nhóm phát triển tập trung vào việc xây dựng các tính năng business logic thay vì quản lý infrastructure.

Kết quả đánh giá và kiểm thử cho thấy hệ thống có thể xử lý ổn định với 50-100 người dùng đồng thời, đạt throughput khoảng 500 requests mỗi giây cho các API queries, và có khả năng tự động mở rộng thông qua ECS auto-scaling. Các metrics về thời gian phản hồi đều nằm trong phạm vi chấp nhận được, với hầu hết các endpoints có response time dưới 300ms. Việc sử dụng hybrid database approach (PostgreSQL + TimescaleDB) đã chứng minh hiệu quả trong việc xử lý cả dữ liệu quan hệ và time-series, với TimescaleDB's continuous aggregates giúp tối ưu hóa các queries trên large time ranges.

Hệ thống đã được kiểm thử về bảo mật và cho thấy không có các vulnerabilities nghiêm trọng. Các biện pháp bảo mật cơ bản như password hashing, session encryption, và secrets management đã được implement đúng cách. Tuy nhiên, một số tính năng bảo mật nâng cao như CSRF protection và rate limiting chưa được implement, và cần được thêm vào trong các phiên bản tiếp theo.

Mặc dù còn một số điểm cần cải thiện về hiệu năng và bảo mật, hệ thống đã chứng minh được tính khả thi và có tiềm năng phát triển lớn trong tương lai. Kiến trúc microservices và việc sử dụng các công nghệ hiện đại tạo nền tảng vững chắc cho việc mở rộng và phát triển thêm các tính năng. Hệ thống đã sẵn sàng để được sử dụng trong môi trường development và testing, và với một số cải thiện, có thể được triển khai trong môi trường production.

\subsection{Khuyến Nghị}

Dựa trên kết quả đánh giá và kiểm thử, các khuyến nghị được chia thành ba categories: khuyến nghị cho phát triển tiếp theo, khuyến nghị cho nghiên cứu tiếp theo, và khuyến nghị cho ứng dụng thực tế. Mỗi category có các priorities khác nhau và có thể được thực hiện song song hoặc tuần tự tùy thuộc vào resources và timeline.

\textbf{Cho Phát Triển Tiếp Theo:}
\begin{itemize}
    \item \textbf{Triển khai lớp cache Redis}: Implement Redis caching layer để cải thiện hiệu năng cho các queries thường xuyên. Redis có thể được sử dụng để cache kết quả của các queries phức tạp, session data, và frequently accessed data. Điều này có thể giảm đáng kể load trên database và cải thiện response times. Redis có thể được deploy như một ECS service hoặc sử dụng AWS ElastiCache.
    
    \item \textbf{Thêm các bài kiểm thử tự động toàn diện}: Thêm comprehensive unit tests và integration tests để đảm bảo chất lượng code và giảm thiểu regressions. Unit tests nên cover tất cả business logic, và integration tests nên cover các workflows chính. Tests nên được chạy tự động trong CI/CD pipeline để catch bugs sớm.
    
    \item \textbf{Triển khai rate limiting và CSRF protection}: Implement rate limiting để bảo vệ hệ thống khỏi abuse và DDoS attacks. CSRF protection nên được thêm vào cho tất cả state-changing operations. Có thể sử dụng middleware hoặc API gateway để implement các tính năng này.
    
    \item \textbf{Thêm tài liệu API với Swagger/OpenAPI}: Thêm API documentation với Swagger/OpenAPI để cải thiện developer experience. Documentation nên bao gồm tất cả endpoints, request/response schemas, và examples. Có thể generate documentation tự động từ code annotations.
    
    \item \textbf{Tích hợp CloudWatch Alarms}: Tích hợp CloudWatch Alarms cho proactive monitoring. Alarms nên được setup cho các metrics quan trọng như CPU utilization, memory usage, error rate, và response time. Alarms có thể trigger notifications qua email hoặc Slack khi thresholds được exceeded.
\end{itemize}

\textbf{Cho Nghiên Cứu Tiếp Theo:}
\begin{itemize}
    \item \textbf{Nghiên cứu các phương pháp tối ưu hóa truy vấn database}: Nghiên cứu các phương pháp tối ưu hóa database queries cho time-series data lớn, bao gồm việc sử dụng indexes hiệu quả, query optimization techniques, và materialized views. So sánh hiệu năng của các approaches khác nhau và identify best practices.
    
    \item \textbf{Đánh giá hiệu quả của các chiến lược cache}: Đánh giá hiệu quả của các caching strategies khác nhau, bao gồm in-memory caching, distributed caching, và CDN caching. So sánh trade-offs giữa cache hit rate, consistency, và complexity.
    
    \item \textbf{So sánh chi phí giữa self-managed và managed services}: So sánh chi phí giữa self-managed databases (PostgreSQL trên ECS) và managed services (AWS RDS). Phân tích trade-offs giữa cost, performance, và operational overhead. Đánh giá khi nào nên sử dụng managed services vs self-managed.
    
    \item \textbf{Nghiên cứu edge computing}: Nghiên cứu edge computing để giảm latency cho real-time applications. Đánh giá khả năng deploy một phần logic xử lý gần với devices để giảm network latency và improve user experience.
    
    \item \textbf{Đánh giá khả năng tích hợp machine learning}: Đánh giá khả năng tích hợp machine learning cho predictive analytics, anomaly detection, và data classification. Nghiên cứu các frameworks và services có sẵn như AWS SageMaker và TensorFlow.
\end{itemize}

\textbf{Cho Ứng Dụng Thực Tế:}
\begin{itemize}
    \item \textbf{Triển khai trên môi trường production}: Triển khai hệ thống trên production environment với proper monitoring, alerting, và backup strategies. Production deployment nên có separate infrastructure từ development/staging, với stricter security policies và higher availability requirements.
    
    \item \textbf{Thực hiện security audit và penetration testing}: Thực hiện security audit và penetration testing bởi các chuyên gia bảo mật bên ngoài để identify và fix các vulnerabilities. Audit nên bao gồm code review, infrastructure review, và penetration testing.
    
    \item \textbf{Tối ưu hóa chi phí hạ tầng đám mây}: Tối ưu hóa chi phí cloud infrastructure bằng cách analyze usage patterns, identify unused resources, và optimize resource allocation. Có thể sử dụng AWS Cost Explorer và Reserved Instances để giảm costs.
    
    \item \textbf{Phát triển ứng dụng di động}: Phát triển mobile applications (iOS và Android) để mở rộng user base và improve accessibility. Mobile apps có thể sử dụng cùng RESTful API, và có thể implement offline capabilities và push notifications.
    
    \item \textbf{Tích hợp với các hệ thống bên thứ ba}: Tích hợp với các hệ thống third-party như weather services, mapping services, và notification services để tăng tính hữu ích của hệ thống. Integrations có thể được implement như microservices riêng biệt để maintain separation of concerns.
\end{itemize}

